\section{Use a direct-mapped cache, not a palette}
After reconstructing a pixel, place its complete RGBA value into slot
\[
h=(3r+5g+7b+11a)\bmod64.
\]
The largest weighted sum is 6630, so the local computation needs no wide multiplication
result. Since the divisor is a power of two, the low six bits select the slot. Multiply
the slot number by four only when forming the byte address of its word. The slot number
itself remains from zero through sixty-three.

An INDEX command uses its low six bits as the slot to fetch. It restores alpha along
with RGB. It does not recompute the original color from that number. Two different
pixels can have the same hash; storing the later one replaces the earlier entry.
A matching hash is therefore insufficient for an encoder. Compare all four components
before selecting INDEX.

For opaque black, the hash is \(11\cdot255\bmod64=53\). Opaque red hashes to fifty.
The cache begins with transparent black even in slot fifty-three; it becomes opaque
black there only after that pixel has been reconstructed. Reassigning the same entry
while expanding a run changes no values, but the cache must have the right state by
the time the following command is decoded.

The specification forbids consecutive INDEX chunks selecting the same slot; RUN is
the representation for that repetition. The strict validator in this exercise checks
that rule. Remember whether the immediately preceding chunk was INDEX and its slot.
A different intervening chunk ends that adjacency. Do not prohibit a later reuse of
the same slot after intervening pixels.
